% ===========================================================================
% Section 5 -- Proposed improvements
% ===========================================================================

We propose two improvements, each aimed at the bottlenecks B1--B5 of
Section~\ref{sec:bottlenecks}:
\begin{itemize}[nosep]
  \item \textbf{(CF+)}, a strengthened compact model (Track A, with the warm
    start of Track C): clique constraints (B3), symmetry breaking (B2), a
    smaller model (B1) and the greedy solution as a starting point (B4);
  \item \textbf{(SP-CG)}, column generation for the set partitioning model
    (Track B), which removes the employee index altogether (B1, B2), has a
    much stronger bound (B3) and is exact on any graph (B5). Optimality is
    proved by \emph{reduced-cost enumeration} instead of a full
    branch-and-price.
\end{itemize}
The code is in \texttt{src/models/improved.py} and
\texttt{src/models/colgen.py}.

\subsection{Strengthened conflict formulation (CF+)}
\label{sec:cfp}

(CF+) has the same variables $y_k$, $x_{ik}$ as (CF) and is exact under the
same condition (transitive $G$). It differs in three ways.

\paragraph{Clique constraints.} Let $\mathcal Q$ be a family of cliques of
the time-conflict graph (duties $i,j$ with neither $(i,j)$ nor $(j,i)$ in
$\A$) such that every conflicting pair lies in some $Q\in\mathcal Q$. The
linking and time-conflict rows \eqref{cf:link}--\eqref{cf:conflict} are
replaced by
\begin{equation}\label{cfp:clique}
  \sum_{i\in Q} x_{ik}\le y_k \qquad Q\in\mathcal Q,\; k\in\K .
\end{equation}
Each row of~\eqref{cfp:clique} implies the pairwise rows of all pairs in $Q$,
and, since every duty lies in some clique, also $x_{ik}\le y_k$; the LP
relaxation is therefore at least as strong as that of (CF), and strictly
stronger as soon as a clique has three or more duties. This is exactly the
weakness of Proposition~\ref{prop:lp-cf}: with $q$ pairwise overlapping
duties, \eqref{cfp:clique} gives $\sum_i x_{ik}\le y_k$, so the fractional
solution of the proof is cut off.

When transfer times are zero, the conflict graph is the overlap graph of the
duty intervals, an interval graph, whose maximal cliques are the sets of
duties running at a common moment, at most $n$ of them. With transfer
times, two duties that do not overlap can still conflict ($e_i<s_j<e_i+t_{ij}$),
and the conflict graph is only a co-comparability graph, which can have
exponentially many maximal cliques. We therefore build $\mathcal Q$ in two
steps: first the interval cliques (the set of duties running at each start
time, keeping only maximal ones), then a greedy cover of the remaining
conflicting pairs, each grown into a clique by adding common neighbours.
On our instances $|\mathcal Q|$ is between $0.6n$ and $1.4n$.

\paragraph{Spread conflicts.} The pairwise spread rows of (CF) are kept.
We also implemented an alternative with $O(n)$ rows per employee: with duties
sorted by start time, a continuous ``started before'' variable
$\sigma_{qk}\ge\sigma_{q-1,k}$, $\sigma_{qk}\ge x_{(q)k}$ and, for each duty
$j$, $x_{jk}+\sigma_{p(j),k}\le y_k$, where $p(j)$ is the last position with
$s<e_j-T_k$. Setting $\sigma_{qk}=\max_{l\le q}x_{(l)k}$ shows that its LP
projection on $(x,y)$ is exactly the set of pairwise spread rows, so the LP
bound is the same with far fewer rows. In practice it was worse: the
pairwise rows are binary conflicts that HiGHS adds to its \emph{clique table},
from which its root cuts combine time and spread conflicts. Hidden behind
$\sigma$, they are lost and the root bound drops (Section~\ref{sec:ablation}).

\paragraph{Symmetry breaking.} Within each type, number the employees
$0,1,\dots$ and the duties by start time. Any solution can be relabelled so
that the employees used come first, ordered by their first duty. Then
employee $r$ has at least $r$ duties starting before its first one, so its
first duty has position at least $r$. Hence the constraints
\begin{equation}\label{cfp:sym}
  x_{ik}=0 \;\;\text{if } \mathrm{pos}(i)<r(k),\qquad
  y_{k}\le y_{k-1}\;\;(r(k)\ge 1)
\end{equation}
remove symmetric copies but keep at least one copy of every solution. The
fixed variables are not even created. On our instances this removes about
30\% of the variables, and the rows shrink by a factor of 2.6, since the
conflict rows of an employee only involve the duties it can still take.

\paragraph{Starting solution: an improving cut.} The greedy solution
$\bar x$ (with cost $\bar z$) is known before the solve. The natural way to
use it is a MIP start, but HiGHS~1.16 handles MIP starts incorrectly: when
presolve fixes variables to $1$, it evaluates the start without the constant
cost of those variables, and then prunes with a wrong incumbent value. On a
set partitioning model it returned an ``optimal'' value of $8.8$ for an
instance whose optimum is $9.6$. We use instead a solver-independent
construction that is correct by design. Every cost is a multiple of $0.2$, so
a solution better than $\bar x$ costs at most $\bar z-0.2$, and we add the
\emph{improving cut}
\begin{equation}\label{cfp:cut}
  \textstyle\sum_k c_k y_k\le \bar z-0.2 .
\end{equation}
If the model with~\eqref{cfp:cut} is infeasible, $\bar x$ is optimal; if the
time runs out without a solution, $\bar x$ is returned with the solver's
bound. In both cases the run always ends with a solution at least as good as
the greedy, and the cut also helps the solver prune.

\paragraph{Size.} With pairwise spread conflicts the model is still
$O(|\K|\,n^2)$ in the worst case, but symmetry breaking and the clique rows
make it 2--7 times smaller than (CF) (median 3.4, in non-zeros). The same
4-million non-zero limit applies.

\subsection{Column generation with reduced-cost enumeration (SP-CG)}
\label{sec:cg}

\paragraph{Master problem.} The restricted master problem (RMP) is the LP
relaxation of (SP) over a subset of columns, with equality rows
$\sum_p a_{ip}\lambda_p=1$. Covering rows ($\ge1$) would allow removing a
duty from a shift, which is only valid if $G$ is transitive, so we keep
equalities and (SP-CG) is exact on every DAG. Each column is a path of $G$
with cost $0.6$ if it is feasible for a half-time employee and $1$ otherwise.
The RMP starts from the $n$ single-duty shifts, which guarantee feasibility,
and the shifts of the greedy solution. It is solved by the simplex method of
HiGHS, warm-started after each addition of columns.

\paragraph{Pricing.} Given the duals $\pi$, for each type $\tau$ we look for
paths $p$ of $G$ with negative reduced cost $c_\tau-\sum_{i\in p}\pi_i$ that
satisfy (S2)--(S3) for $\tau$. Labels $(\text{cost}=-\sum\pi,\;\text{work},\;
s_{\mathrm f},\;\text{path})$ are extended along the arcs of $G$ in order of
start time (a topological order), the standard labelling scheme for this
problem~\cite{irnich2005}. Three devices make this fast enough:
\begin{itemize}[nosep]
  \item \emph{Grouping by first start.} The spread of every extension
    depends only on $s_{\mathrm f}$ and the end node, so labels at a node are
    grouped by $s_{\mathrm f}$. In each group we keep the Pareto front in
    (work, cost), stored sorted by work, so a new label is inserted or
    rejected by binary search. A label is discarded only if a label with the
    same $s_{\mathrm f}$ has no larger work and cost, so the pricing is exact.
  \item \emph{Completion bound.} Let
    $g(j)=\min\{0,\min_{(j,k)\in\A}(g(k)-\pi_k)\}$, computed backwards in
    $O(|\A|)$. It is the cheapest cost of any extension of a label at $j$,
    resources ignored. A label with $c_\tau+\text{cost}+g(j)\ge0$ cannot lead
    to a negative column and is dropped.
  \item \emph{Heuristic pass.} First only the cheapest label per group, and
    the 30 cheapest groups per node, are extended. The exact pass runs only
    when this finds no negative column, so the last iteration is exact.
\end{itemize}
Up to 300 columns per type are added in each iteration. When the exact
pricing finds no negative column, the RMP value is $z_{\mathrm{LP}}$, the LP
bound of (SP). Before that, the exact minimum reduced costs
$\bar c_\tau$ give the Lagrangian bound~\cite{lubbecke2005}
$z_{\mathrm{RMP}}+\sum_\tau(\bar z/c_\tau)\min(0,\bar c_\tau)$, since a
solution of cost at most $\bar z$ uses at most $\bar z/c_\tau$ shifts of
type $\tau$. Since all costs are multiples of $0.2$, the bound used is
$\lceil 5z_{\mathrm{LP}}\rceil/5$.

\paragraph{Integer solutions: price-and-branch.} The RMP with all the
generated columns is solved as a MILP (with the improving
cut~\eqref{cfp:cut} for the greedy solution), using at most 40\% of the
remaining time. This gives a good upper bound $\bar z$ but, in general, no
proof of optimality: an optimal solution may need columns that were never
generated.

\paragraph{Optimality proof: reduced-cost enumeration.} Instead of
branch-and-price, we use the following argument, which underlies the
column-enumeration methods of exact set partitioning
algorithms~\cite{baldacci2008}.

\begin{proposition}\label{prop:enum}
Let $\pi$ be optimal duals of the LP relaxation of (SP), with value
$z_{\mathrm{LP}}=\sum_i\pi_i$, and let $\bar c_p=c_p-\sum_{i\in p}\pi_i\ge0$
be the reduced costs. Every solution of (SP) with cost at most $\bar z-0.2$
uses only columns with $\bar c_p\le\Delta:=\bar z-0.2-z_{\mathrm{LP}}$.
\end{proposition}
\begin{proof}
Let $P$ be the set of columns of a solution. Since $P$ partitions $\D$,
$\sum_{p\in P}c_p=\sum_{p\in P}\bar c_p+\sum_{p\in P}\sum_{i\in p}\pi_i
=\sum_{p\in P}\bar c_p+z_{\mathrm{LP}}$. If the cost is at most $\bar z-0.2$,
then $\sum_{p\in P}\bar c_p\le\Delta$, and each term is non-negative, so each
is at most $\Delta$.
\end{proof}

So, after column generation, we enumerate \emph{all} feasible paths with
$\bar c_p\le\Delta$ by depth-first search from every duty, pruned with the
completion bound $g$ (a partial path is abandoned when
$c_{\min}+\text{cost}+g(j)>\Delta$, where $c_{\min}=0.6$ while the path is
still feasible for a half-time employee and $1$ afterwards). The MILP over
these columns is then solved. If it finds a solution better than $\bar z$,
that solution is optimal; if it has none, $\bar z$ is optimal. The
enumeration is small when the gap $\bar z-z_{\mathrm{LP}}$ is small, which is
the case here because the LP bound of (SP) is very tight
(Section~\ref{sec:comparison}). If more than $10^6$ columns qualify, or the
time limit is reached, we stop and report the gap between $\bar z$ and
$\lceil 5z_{\mathrm{LP}}\rceil/5$.

\paragraph{Exactness.} The procedure is exact on every DAG: the RMP has the
partitioning rows, the pricing and the enumeration only discard labels that
are dominated or cannot meet the threshold, and the final claim of
optimality rests on Proposition~\ref{prop:enum}. Unlike the compact models it
needs no employee pool, so it has no symmetry and its size does not depend on
$|\K|$.

\subsection{Verification of the improvements}
\label{sec:verif-improved}

The test script of Section~\ref{sec:verification} was extended to the new
models. On the nine hand-made instances, (SP-CG) returns the exact optimum,
including $0.6$ on the non-transitive one, and every variant of (CF+) returns
the value of (CF). In addition, eight generated instances ($n=12$--$26$,
three of them non-transitive with arc\_drop $=0.3$) are solved by all
methods: (SP-CG) always matches (AF), and all (CF+) variants match (CF). In
the full experiment, every solution is checked by the validator, and the
analysis checks that all proven optima of an instance agree
(Section~\ref{sec:comparison}).
